\clearpage
\phantomsection% 
\addcontentsline{toc}{chapter}{ABSTRAK}
\thispagestyle{fancy}

\begin{center}
	\large \bfseries \MakeUppercase{Abstrak}\\
%	\normalsize \normalfont {\thetitle}\\
%	\normalsize \normalfont {\theauthor}\\
	\bigskip
	
	\normalsize \normalfont \justifying \singlespacing
	Penerimaan Mahasiswa Baru (PMB) merupakan proses krusial dalam pengelolaan perguruan tinggi yang memerlukan sistem terstruktur dan efisien. Universitas HKBP Nommensen saat ini masih mengandalkan proses manual dalam pelaksanaan PMB, yang menyebabkan berbagai kendala seperti lambatnya pemrosesan data pendaftar, kesulitan dalam verifikasi berkas, serta kurangnya transparansi informasi bagi calon mahasiswa. Penelitian ini bertujuan untuk mengembangkan sistem terintegrasi penerimaan mahasiswa baru berbasis web yang mampu mengelola seluruh proses PMB secara digital, mulai dari pendaftaran, verifikasi data, pembayaran, hingga daftar ulang. Metode pengembangan yang digunakan adalah metode \textit{Prototype}, yang dipilih karena memungkinkan pengembangan bertahap dengan umpan balik langsung dari pengguna. Sistem dirancang menggunakan arsitektur \textit{Model-View-Controller} (MVC) dengan \textit{framework} Spring Boot, basis data MySQL, serta antarmuka \textit{responsive} menggunakan HTML, CSS, dan Bootstrap. Perancangan sistem meliputi analisis kebutuhan fungsional dan non-fungsional, pembuatan diagram \textit{use case}, \textit{activity diagram}, \textit{sequence diagram}, \textit{class diagram}, \textit{Data Flow Diagram} (DFD), \textit{Entity Relationship Diagram} (ERD), serta desain antarmuka menggunakan Figma. Pengujian sistem dilakukan melalui tiga pendekatan, yaitu \textit{Black-box Testing} untuk memvalidasi fungsionalitas, \textit{Integration Testing} untuk menguji keterpaduan antar modul, serta \textit{System Usability Scale} (SUS) untuk mengukur tingkat kebergunaan sistem dari perspektif pengguna. Hasil penelitian diharapkan dapat menghasilkan sistem PMB yang efisien, terintegrasi, dan mudah digunakan oleh seluruh pemangku kepentingan di Universitas HKBP Nommensen.\par

    \vspace{20pt}
	\raggedright \textbf{Kata Kunci: penerimaan mahasiswa baru, sistem terintegrasi, \textit{prototype}, Spring Boot, web}
	
	\vfill
	
\end{center}
\clearpage
